\chapter{Vitamin B12 (Cobalamin) Deficiency}
\index{Vitamin B12 (Cobalamin) Deficiency}
\label{sec:Vitamin B12 Deficiency}

\section{Overview}
Vitamin B12 (cobalamin) deficiency is a vitamin deficiency disorder common in the elderly due to atrophic gastritis and declining intrinsic factor production (food-cobalamin malabsorption). Pernicious anemia is an autoimmune disorder with antibodies against gastric parietal cells and/or intrinsic factor, leading to profound B12 deficiency. Other causes include gastric or ileal removal, Crohn disease, bacterial overgrowth, fish tapeworm (Diphyllobothrium latum) infection, strict vegetarian/vegan diet, and metformin use. This metabolic disorder involves dysregulation of normal bodily processes, often requiring ongoing monitoring and management. It is increasingly common due to lifestyle and dietary factors. Metabolic disorders involve disruption of normal biochemical processes essential for health. These conditions may result from enzyme deficiencies, hormonal imbalances, or lifestyle factors. Many metabolic disorders are chronic and require ongoing management. Lifestyle modifications are often the cornerstone of treatment.
\section{Causes}
This condition happens because of deficiency of vitamin B12, which is essential for DNA synthesis, myelin formation, and homocysteine metabolism. Metabolic disorders arise from dysregulation of normal biochemical pathways. This may result from enzyme deficiencies, hormonal imbalances, or lifestyle factors that overwhelm normal homeostatic mechanisms. The underlying mechanisms involve complex interactions between genetic predisposition and environmental factors. Disruption of normal physiological processes leads to the characteristic features of the condition. Multiple contributing factors may act synergistically to trigger or worsen the condition.
\section{Risk Factors}
\begin{itemize}[noitemsep]
  \item Genetic predisposition and family history
  \item Advancing age
  \item Lifestyle factors (diet, exercise, smoking, alcohol)
  \item Environmental exposures
  \item Underlying medical conditions
  \item Medication use
\end{itemize}
\section{Signs and Symptoms}

\subsection{Common symptoms}
\begin{itemize}[noitemsep]
  \item Pain or discomfort in the affected area
  \end{itemize}

\subsection{Emergency warning signs}
\begin{itemize}[noitemsep]
  \item Severe or worsening symptoms of vitamin b12 (cobalamin) deficiency require immediate medical evaluation
\end{itemize}
\section{Progression and Stages}
You may experience megaloblastic anemia and the hallmark distinguishing feature from folate deficiency: nervous system involvement (subacute combined deterioration of the spinal cord manifesting as symmetric paresthesias, loss of vibration and proprioception, loss of coordination, spastic paraparesis, and a positive Romberg sign), peripheral neuropathy, visual disturbances (optic shrinkage or wasting), memory loss, and psychiatric symptoms. The condition may progress through different stages of severity over time. Early stages often involve mild or subtle symptoms that may be overlooked. Moderate stages involve more noticeable symptoms affecting daily function. Advanced stages may involve significant impairment and complications.
\section{Complications}
\begin{itemize}[noitemsep]
  \item Disease progression leading to worsening symptoms
  \item Secondary infections or injuries
  \item Side effects of treatment
  \item Reduced quality of life
  \item Emotional and psychological impact
\end{itemize}
\section{Diagnosis}
\begin{itemize}[noitemsep]
  \item Doctors diagnose this using low serum B12 (less than 200 pg/mL) with elevated serum methylmalonic acid (MMA) and homocysteine
  \item Elevated MMA is the most sensitive indicator.
  \item Comprehensive medical history and physical examination
  \item Laboratory tests to assess organ function
  \item Imaging studies to visualize affected structures
  \item Specialized testing depending on the affected system
  \item Consultation with relevant specialists
\end{itemize}
\section{Prevention}
\begin{itemize}[noitemsep]
  \item Maintain a healthy lifestyle with balanced diet and exercise
  \item Avoid known risk factors
  \item Attend regular medical check-ups for early detection
  \item Follow recommended screening guidelines
\end{itemize}
\section{Treatment Options}
\begin{itemize}[noitemsep]
  \item Treatment focuses on oral B12 1000--2000 \textmu g daily for deficiency without nervous system involvement
  \item For pernicious anemia, nervous system symptoms, or severe deficiency, parenteral B12 (hydroxocobalamin or cyanocobalamin) 1000 \textmu g IM daily for 1 week, weekly for 4 weeks, then monthly for life.
  \item Medications to manage symptoms and address underlying mechanisms
  \item Lifestyle modifications and self-care measures
  \item Physical therapy and rehabilitation
  \item Surgical intervention when indicated
  \item Regular monitoring and follow-up care
\end{itemize}
\section{Living with It}
\begin{itemize}[noitemsep]
  \item Follow your treatment plan as prescribed
  \item Maintain regular follow-up with your healthcare provider
  \item Make necessary lifestyle adjustments
  \item Seek support from family, friends, or support groups
  \item Stay informed about your condition
\end{itemize}
\section{Outlook}
\begin{itemize}[noitemsep]
  \item Most people recover well with treatment
  \item Hematologic response is rapid (reticulocytosis within 3--5 days, hemoglobin normalization in 6--8 weeks).
\end{itemize}
